\section{Related Work in Music Programming and Representation}\label{sec:related-work}

Existing approaches to music programming differ along three axes that matter for this thesis: whether the system
targets live performance or offline artifacts, whether composition is separated from synthesis, and whether temporal
structure is expressed explicitly or through implicit sequencing.
This section surveys representative systems in each category.
Section~\ref{sec:motivo-position} synthesizes the resulting research gap.

\subsection{Live-Coding and Performance-Oriented Languages}\label{subsec:live-coding-related}

\textbf{Sonic Pi} is a Ruby-embedded language designed for education and live performance~\cite{Aaron_Blackwell_2013,
Aaron_2016}.
It wraps the SuperCollider synthesis server and exposes a sequential execution model: \texttt{play} triggers a note,
\texttt{sleep} advances time, and simple loops provide repetition.
The design prioritizes accessibility for beginners and low-latency classroom use on constrained hardware.
Sonic Pi is effective as a pedagogical live-coding environment, but its sequential model and dependence on an external
synthesizer make it a weak fit for offline structural composition with a portable MIDI artifact as the primary output.

\textbf{TidalCycles} (often called Tidal) is a Haskell-embedded pattern language for live
coding~\cite{McLean_Wiggins_2010}.
Musical structures are represented as time-varying patterns that can be transformed, concatenated, and combined through
functional composition.
Tidal's random-access pattern functions avoid the recomputation cost of lazy lists during performance, but the language
streams events to external synthesizers via Open Sound Control rather than compiling to a self-contained file format.
\textbf{Strudel} ports Tidal's pattern model to JavaScript and the browser, supporting WebAudio, OSC, and optional MIDI
output among several backends~\cite{Roos_McLean_2023}.
Both Tidal and Strudel remain oriented toward live manipulation; MIDI, when used, is one of several runtime outputs
rather than the deterministic compilation target.

\textbf{HMusic} embeds a grid-like pattern and track model in Haskell, resembling sequencer interfaces while retaining
functional abstractions~\cite{Du_Bois_Ribeiro_2019}.
Patterns define rhythmic hits and rests; tracks bind patterns to instruments.
The language compiles to Sonic Pi for playback, which again couples composition to a live synthesis pipeline rather
than to a standalone symbolic artifact.

\textbf{ChucK} is a strongly timed, concurrent language for real-time audio synthesis~\cite{Wang_Cook_2003,Wang_2008}.
Time is a first-class value advanced explicitly through the \texttt{now} keyword; lightweight processes called
\textit{shreds} run concurrently with deterministic synchronization.
ChucK excels at precise on-the-fly audio programming but operates at a low level of abstraction: the programmer
manages time advancement and signal routing directly, which makes large-scale structural composition harder to express
than short synthesis patches.

\textbf{SuperCollider} combines a high-level language with a separate synthesis server communicating over
OSC~\cite{McCartney_2002,Wilson_Cottle_Collins_2011}.
Unit generators compose into synthesis graphs; higher-level \textit{patterns} and \textit{streams} generate event
sequences over time.
SuperCollider is widely used for algorithmic composition and live performance, yet its primary abstractions remain
signal-processing oriented.
Representing hierarchical song structure as an explicit, compilable timeline is not the central design goal.

Across these systems, a common pattern emerges: composition and execution are intertwined.
The composer's program drives a running audio engine or event stream.
Motivo instead compiles source to a finished MIDI file with no runtime dependency on Motivo itself.

\subsection{Notation and Interchange Formats}\label{subsec:notation-related}

\textbf{ABC notation} is a plain-text format for notating folk and traditional music~\cite{Walshaw_2011}.
Header fields encode metadata (title, meter, key); the body encodes notes, durations, and repeats using an ASCII
alphabet.
ABC is human-readable and widely supported by conversion tools, but it is a notation format rather than a programming
language: it lacks parameterized patterns, typed variables, and imperative control flow.
MusicXML addresses score interchange between notation programs through a structured XML representation~\cite{Good_2001}.

As discussed in Section~\ref{sec:midi-background}, \textbf{MIDI} encodes performance events with strong
interoperability~\cite{Moog_1986}.
Moore argues that MIDI conflates musical structure with performance data and lacks native support for higher-level
abstractions such as phrasing and hierarchy~\cite{Moore_1988}.
MIDI is therefore an excellent \textit{backend} target for rendering and interchange, but a poor \textit{primary}
representation for structured algorithmic composition.
Motivo uses MIDI as its output format while keeping compositional intent in the source language.
